\documentclass[12pt,a4paper]{ctexbook}
\usepackage{geometry}
\geometry{margin=2.2cm}
\usepackage{setspace}
\setstretch{1.35}
\usepackage{hyperref}
\title{千家诗}
\author{古典文献整理}
\date{}
\begin{document}
\maketitle
\tableofcontents
\newpage
\section{第1则}
春眠·（唐）孟浩然\par
春眠不覺曉，處處聞啼鳥。\par
夜來風雨聲，花落知多少。\par
訪袁拾遺不遇·（唐）孟浩然\par
洛陽訪才子，江嶺作流人。\par
聞說梅花早，何如此地春。\par
送郭司倉·（唐）王昌齡\par
映門淮水綠，留騎主人心。\par
明月隨良掾，春潮夜夜深。\par
洛陽道·（唐）儲光羲\par
大道直如發，春來佳氣多。\par
五陵貴公子，雙雙嗚玉珂。\par
獨坐敬亭山·（唐）李白\par
衆鳥高飛盡，孤雲獨去閒。\par
相看兩不厭，只有敬亭山。\par
登鸛雀樓·（唐）王之渙\par
白日依山盡，黃河入海流。\par
欲窮千里目，更上一層樓。\par
觀永樂公主入番·（唐）孫逖\par
邊地鶯花少，年來未覺新。\par
美人天上落，龍塞始應春。\par
伊州歌·（唐）金昌緒\par
打起黃鶯兒，莫教枝上啼。\par
啼時驚妾夢，不得到遼西。\par
左掖梨花·（唐）丘爲\par
冷豔全欺雪，餘香乍入衣。\par
春風且莫定，吹向玉階飛。\par
思君恩·（唐）令狐楚\par
小苑鶯歌歇，長門蝶舞多。\par
眼看春又去，翠輦不曾過。\par
題袁氏別業·（唐）賀知章\par
主人不相識，偶坐爲林泉。\par
莫謾愁沽酒，囊中自有錢。\par
夜送趙縱·（唐）楊炯\par
趙氏連城璧，由來天下傳。\par
送君還舊府，明月滿前川。\par
竹裏館·（唐）王維\par
獨坐幽篁裏，彈琴復長嘯。\par
深林人不知，明月來相照。\par
送朱大入秦·（唐）王維\par
避人五陵去，寶劍值千金。\par
分手脫相贈，平生一片心。\par
長幹行·（唐）崔顥\par
君家在何處？妾住在橫塘。\par
停船暫借問，或恐是同鄉。\par
詠史·（唐）高適\par
尚有綈袍贈，應憐範叔寒。\par
不知天下士，猶作布衣看。\par
罷相作·（唐）李適之\par
避賢初罷相，樂聖且銜杯。\par
爲問門前客，今朝幾個來。\par
逢俠者·（唐）錢起\par
燕趙悲歌士，相逢劇孟家。\par
寸心言不盡，前路日將斜。\par
江行望匡廬·（唐）錢起\par
咫尺愁風雨，匡廬不可登。\par
只疑雲霧窟，猶有六朝僧。\par
答李浣·（唐）韋應物\par
林中觀易罷，溪上對鷗閒。\par
楚俗饒詞客，何人最往還。\par
秋風引·（唐）劉禹錫\par
何處秋風至，蕭蕭送雁羣。\par
朝來入庭樹，孤客最先聞。\par
秋夜寄丘員外·（唐）韋應物\par
懷君屬秋夜，散步詠涼天。\par
山空松子落，幽人應未眠。\par
秋日·（唐）耿湋\par
返照入閭巷，憂來誰共語。\par
古道少人行，秋風動禾黍。\par
秋日湖上·（唐）薛瑩\par
落日五湖遊，煙波處處愁。\par
浮沉千古事，誰與問東流。\par
宮中題·（唐）李昂\par
輦路生秋草，上林花滿枝。\par
憑高何限意，無復侍臣知。\par
汾上驚秋·（唐）蘇頲\par
北風吹白雲，萬里渡河汾。\par
心緒逢搖落，秋聲不可聞。\par
尋隱者不遇·（唐）賈島\par
松下問童子，言師採藥去。\par
只在此山中，雲深不知處。\par
蜀道後期·（唐）張說\par
客心爭日月，來往預期程。\par
秋風不相待，先至洛陽城。\par
靜夜思·（唐）李白\par
牀前明月光，疑是地上霜。\par
舉頭望明月，低頭思故鄉。\par
秋浦歌·（唐）李白\par
白髮三千丈，離愁似個長。\par
不知明鏡裏，何處得秋霜。\par
贈喬侍郎·（唐）陳子昂\par
漢廷榮巧宦，雲閣薄邊功。\par
可憐驄馬使，白首爲誰雄。\par
答五陵太守·（唐）王昌齡\par
仗劍行千里，微軀敢一言。\par
曾爲大梁客，不負信陵恩。\par
行軍九日思長安故園·（唐）岑參\par
強欲登高去，無人送酒來。\par
遙憐故園菊，應傍戰場開。\par
婕妤怨·（唐）皇甫冉\par
花枝出建章，鳳管發昭陽。\par
借問承恩者，雙蛾幾許長？\par
題竹林寺·（唐）朱放\par
歲月人間促，煙霞此地多。\par
殷勤竹林寺，更得幾回過？\par
過三閭廟·（唐）戴叔倫\par
沅湘流不盡，屈子怨何深！\par
日暮秋風起，蕭蕭楓樹林。\par
易水送別·（唐）駱賓王\par
此地別燕丹，壯士發衝冠。\par
昔時人已沒，今日水猶寒。\par
別盧秦卿·（唐）司空曙\par
知有前期在，難分此夜中。\par
無將故人酒，不及石尤風。\par
答人·（唐）太上隱者\par
偶來松樹下，高枕石頭眠。\par
山中無曆日，寒盡不知年。\par

\section{第2则}
幸蜀回至劍門·（唐）李隆基\par
劍閣橫雲峻，鑾輿出狩回。\par
翠屏千仞合，丹嶂五丁開。\par
灌木縈旗轉，仙雲拂馬來。\par
乘時方在德，嗟爾勒銘才。\par
和晉陵陸承相·（唐）杜審言\par
獨有宦遊人，偏驚物候新。\par
雲霞出海曙，梅柳渡江春。\par
淑氣催黃鳥，晴光轉綠蘋。\par
忽聞歌古調，歸思欲沾巾。\par
蓬萊三殿侍宴奉敕詠終南山·（唐）杜審言\par
北斗掛城邊，南山倚殿前。\par
雲標金闕回，樹杪玉堂懸。\par
半嶺通佳氣，中峯繞瑞煙。\par
小臣持獻壽，長此戴堯天。\par
春夜別友人·（唐）陳子昂\par
銀燭吐青煙，金尊對綺筵。\par
離堂思琴瑟，別路繞山川。\par
明月懸高樹，長河沒曉天。\par
悠悠洛陽道，此會在何年。\par
長寧公主東莊侍宴·（唐）李嶠\par
別業臨青甸，鳴鑾降紫霄。\par
長筵鵷鷺集，仙管鳳凰調。\par
樹接南山近，煙含北渚遙。\par
承恩鹹已醉，戀賞未還鑣。\par
恩賜麗正殿書院賜宴應制得林字·（唐）張說\par
東壁圖書府，西園翰墨林。\par
誦詩聞國政，講易見天心。\par
位竊和羹重，恩叨醉酒深。\par
載歌春興曲，情竭爲知音。\par
送友人·（唐）李白\par
青山橫北郭，白水繞東城。\par
此地一爲別，孤篷萬里徵。\par
浮雲遊子意，落日故人情。\par
揮手自茲去，蕭蕭斑馬鳴。\par
送友人入蜀·（唐）李白\par
見說蠶叢路，崎嶇不易行。\par
山從人面起，雲傍馬頭生。\par
芳樹籠秦棧，春流繞蜀城。\par
升沈應已定，不必問君平。\par
次北固山下·（唐）王灣\par
客路青山外，行舟綠水前。\par
潮平兩岸闊，風正一帆懸。\par
海日生殘夜，江春入舊年。\par
鄉書何處達，歸雁洛陽邊。\par
蘇氏別業·（唐）祖詠\par
別業居幽處，到來生隱心。\par
南山當戶牖，澧水映園林。\par
竹覆經冬雪，庭昏未夕陰。\par
寥寥人境外，閒坐聽春禽。\par
春宿左省·（唐）杜甫\par
花隱掖垣暮，啾啾棲鳥過。\par
星臨萬戶動，月傍九霄多。\par
不寢聽金鑰，因風想玉坷。\par
明朝有封事，數問夜如何？\par
題玄武禪師屋壁·（唐）杜甫\par
何年顧虎頭，滿壁畫滄州。\par
赤日石林氣，青天江海流。\par
錫飛常近鶴，杯渡不驚鷗。\par
似得廬山路，真隨惠遠遊。\par
終南山·（唐）王維\par
太乙近天都，連山到海隅。\par
白雲回望合，青靄入看無。\par
分野中峯變，陰晴衆壑殊。\par
欲投何處宿，隔水問樵夫。\par
登總持閣·（唐）岑參\par
高閣逼諸天，登臨近日邊。\par
晴開萬井樹，愁看五陵煙。\par
檻外低秦嶺，窗中小渭川。\par
早知清淨理，常願奉金仙。\par
寄左省杜拾遺·（唐）岑參\par
聯步趨丹陛，分曹限紫薇。\par
曉隨天仗入，暮惹御香歸。\par
白髮悲花落，青雲羨鳥飛。\par
聖朝無闕事，自覺諫書稀。\par
登兗州城樓·（唐）杜甫\par
東郡趨庭日，南樓縱目初。\par
浮雲連海岱，平野入青徐。\par
孤嶂秦碑在，荒城魯殿餘。\par
從來多古意，臨眺獨躊躇。\par
杜少府之任蜀州·（唐）王勃\par
城闕輔三秦，風煙望五津。\par
與君離別意，同是宦遊人。\par
海內存知己，天涯若比鄰。\par
無爲在歧路，兒女共沾巾。\par
送崔融·（唐）杜審言\par
君王行出將，書記遠從徵。\par
祖帳連河闕，軍麾動洛城。\par
旌旗朝朔氣，笳吹夜邊聲。\par
坐覺煙塵少，秋風古北平。\par
扈從登封途中作·（唐）宋之問\par
帳殿鬱崔嵬，仙遊實壯哉！\par
曉雲連幕卷，夜火雜星迴。\par
谷暗千旗出，山鳴萬乘來。\par
扈從良可賦，終乏掞天才。\par
題義公禪房·（唐）孟浩然\par
義公習禪寂，結宇依空林。\par
戶外一峯秀，階前衆壑深。\par
夕陽連雨足，空翠落庭陰。\par
看取蓮花淨，方知不染心。\par
醉後贈張九旭·（唐）高適\par
世上漫相識，此翁殊不然。\par
興來書自聖，醉後語尤顛。\par
白髮老閒事，青雲在目前。\par
牀頭一壺酒，能更幾回眠。\par
玉臺觀·（唐）高適\par
浩劫因王造，平臺訪古遊。\par
彩雲蕭史駐，文字魯恭留。\par
宮闕通羣帝，乾坤到十洲。\par
人傳有笙鶴，時過北山頭。\par
觀李固請司馬弟山水圖·（唐）杜甫\par
方丈渾連水，天台總映雲。\par
人間長見畫，老去限空聞。\par
范蠡舟偏小，王喬鶴不羣。\par
此生隨萬物，何處出塵氛。\par
旅夜書懷·（唐）杜甫\par
細草微風岸，危檣獨夜舟。\par
星垂平野闊，月涌大江流。\par
名豈文章著，官因老病休。\par
飄飄何所似，天地一沙鷗。\par
登岳陽樓·（唐）杜甫\par
昔聞洞庭水，今上岳陽樓。\par
吳楚東南坼，乾坤日月浮。\par
親朋無一字，老病有孤舟。\par
戎馬關山北，憑軒涕泗流。\par
江南旅情·（唐）祖詠\par
楚山不可極，歸路但蕭條。\par
海色晴看雨，江聲夜聽潮。\par
劍留南鬥近，書寄北風遙。\par
爲報空潭橘，無媒寄洛橋。\par
宿龍興寺·（唐）綦毋潛\par
香剎夜忘歸，鬆清古殿扉。\par
燈明方丈室，珠系比丘衣。\par
白日傳心淨，青蓮喻法微。\par
天花落不盡，處處鳥銜飛。\par
破山寺後禪院·（唐）常建\par
清晨入古寺，初日照高林。\par
曲徑通幽處，禪房花木深。\par
山光悅鳥性，潭影空人心。\par
萬籟此俱寂，惟聞鐘磬音。\par
題鬆汀驛·（唐）張佑\par
山色遠含空，蒼茫澤國東。\par
海明先見日，江白迥聞風。\par
鳥道高原去，人煙小徑通。\par
那知舊遺逸，不在五湖中。\par
聖果寺·（唐）釋處默\par
路自中峯上，盤迴出薜蘿。\par
到江吳地盡，隔岸越山多。\par
古木叢青藹，遙天浸白波。\par
下方城郭近，鐘磬雜笙歌。\par
野望·（唐）王績\par
東皋薄暮望，徙倚欲何依。\par
樹樹皆秋色，山山惟落暉。\par
牧人驅犢返，獵馬帶禽歸。\par
相顧無相識，長歌懷采薇。\par
送著作佐郎崔融等從樑王東征·（唐）陳子昂\par
金天方肅殺，白露始專征。\par
王師非樂戰，之子慎佳兵。\par
海氣侵南部，邊風掃北平。\par
莫賣盧龍塞，歸邀麟閣名。\par
攜妓納涼晚際遇雨·（唐）杜甫\par
落日放船好，輕風生浪遲。\par
竹深留客處，荷淨納涼時。\par
公子調冰水，佳人雪藕絲。\par
片雲頭上黑，應是雨催詩。\par
雨來沾席上，風急打船頭。\par
越女紅裙溼，燕姬翠黛愁。\par
纜侵堤柳系，幔捲浪花浮。\par
歸路翻蕭颯，陂塘五月秋。\par
宿雲門寺閣·（唐）孫逖\par
香閣東山下，煙花象外幽。\par
懸燈千嶂夕，卷幔五湖秋。\par
畫壁餘鴻雁，紗窗宿鬥牛。\par
更疑天路近，夢與白雲遊。\par
秋登宣城謝眺北樓·（唐）李白\par
江城如畫裏，山晚望晴空。\par
兩水夾明鏡，雙橋落彩虹。\par
人煙寒橘柚，秋色老梧桐。\par
誰念北樓上，臨風懷謝公。\par
臨洞庭·（唐）孟浩然\par
八月湖水平，涵虛混太清。\par
氣蒸雲夢澤，波撼岳陽城。\par
欲濟無舟楫，端居恥聖明。\par
坐觀垂釣者，徒有羨魚情。\par
過香積寺·（唐）王維\par
不知香積寺，數裏入雲峯。\par
古木無人徑，深山何處鐘。\par
泉聲咽危石，日色冷青松。\par
薄暮空潭曲，安禪製毒龍。\par
送鄭侍御謫閩中·（唐）高適\par
謫去君無恨，閩中我舊過。\par
大都秋雁少，只是夜猿多。\par
東路雲山合，南天瘴癘和。\par
自當逢雨露，行矣順風波。\par
秦州雜詩·（唐）杜甫\par
鳳林戈未息，魚海路常難。\par
候火雲峯峻，懸軍幕井乾。\par
風連西極動，月過北庭寒。\par
故老思飛將，何時議築壇。\par
禹廟·（唐）杜甫\par
禹廟空山裏，秋風落日斜。\par
荒庭垂橘柚，古屋畫龍蛇。\par
雲氣生虛壁，江深走白沙。\par
早知乘四載，疏鑿控三巴。\par
望秦川·（唐）李頎\par
秦川朝望迥，日出正東峯。\par
遠近山河淨，逶迤城闕重。\par
秋聲萬戶竹，寒色五陵鬆。\par
有客歸歟嘆，悽其霜露濃。\par
同王徵君洞庭有懷·（唐）張謂\par
八月洞庭秋，瀟湘水北流。\par
還家萬里夢，爲客五更愁。\par
不用開書帙，偏宜上酒樓。\par
故人京洛滿，何日復同遊。\par
渡揚子江·（唐）丁仙芝\par
桂楫中流望，空波兩岸明。\par
林開揚子驛，山出潤州城。\par
海盡邊陰靜，江寒朔吹生。\par
更聞楓葉下，淅瀝度秋聲。\par
幽州夜歌·（唐）張說\par
涼風吹夜雨，蕭瑟動寒林。\par
正有高堂宴，能忘遲暮心。\par
軍中宜劍舞，塞上重笳音。\par
不作邊城將，誰知恩遇深。\par

\section{第3则}
春日偶成·（宋）程顥\par
雲淡風輕近午天，傍花隨柳過前川。\par
時人不識餘心樂，將謂偷閒學少年。\par
春日·（宋）朱熹\par
勝日尋芳泗水濱，無邊光景一時新。\par
等閒識得東風面，萬紫千紅總是春。\par
春宵·（宋）蘇軾\par
春宵一刻值千金，花有清香月有陰。\par
歌管樓臺聲細細，鞦韆院落夜沈沈。\par
城東早春·（唐）楊巨源\par
詩家清景在新春，綠柳才黃半未勻。\par
若待上林花似錦，出門俱是看花人。\par
春夜·（宋）王安石\par
金爐香燼漏聲殘，剪剪輕風陣陣寒。\par
春色惱人眠不得，月移花影上欄杆。\par
初春小雨·（宋）韓愈\par
天街小雨潤如酥，草色遙看近卻無。\par
最是一年春好處，絕勝煙柳滿皇都。\par
元日·（宋）王安石\par
爆竹聲中一歲除，春風送暖入屠蘇。\par
千門萬戶曈曈日，總把新桃換舊符。\par
上元侍宴·（宋）蘇軾\par
淡月疏星繞建章，仙風吹下御爐香。\par
侍臣鵠立通明殿，一朵紅雲捧玉皇。\par
立春偶成·（宋）張栻\par
律回歲晚冰霜少，春到人間草木知。\par
便覺眼前生意滿，東風吹水綠參差。\par
打球圖·（宋）晁無咎\par
閶闔千門萬戶開，三郎沈醉打球回。\par
九齡已老韓休死，無復明朝諫疏來。\par
宮詞·（宋）林洪\par
金殿當頭紫閣重，仙人掌上玉芙蓉。\par
太平天子朝元日，五色雲車駕六龍。\par
殿上袞衣明日月，硯中旗影動龍蛇。\par
縱橫禮樂三千字，獨對丹墀日未斜。\par
詠華清宮·（唐）王建\par
行盡江南數十程，曉風殘月入華清。\par
朝元閣上西風急，都入長楊作雨聲。\par
清平調詞·（唐）李白\par
雲想衣裳花想容，春風拂檻露華濃。\par
若非羣玉山頭見，會向瑤臺月下逢。\par
題邸間壁·（宋）鄭會\par
酴醿香夢怯春寒，翠掩重門燕子閒。\par
敲斷玉釵紅燭冷，計程應說到常山。\par
絕句·（唐）杜甫\par
兩個黃鸝鳴翠柳，一行白鷺上青天。\par
窗含西嶺千秋雪，門泊東吳萬里船。\par
海棠·（宋）蘇軾\par
東風嫋嫋泛崇光，香霧空濛月轉廊。\par
只恐夜深花睡去，故燒高燭照紅妝。\par
清明·（宋）王禹偁\par
無花無酒過清明，興味蕭然似野僧。\par
昨日鄰家乞新火，曉窗分與讀書燈。\par
清明·（唐）杜牧\par
清明時節雨紛紛，路上行人慾斷魂。\par
借問酒家何處有，牧童遙指杏花村。\par
社日·（唐）張演\par
鵝湖山下稻粱肥，豚柵雞棲對掩扉。\par
桑柘影斜春社散，家家扶得醉人歸。\par
寒食·（唐）韓翃\par
春城無處不飛花，寒食東風御柳斜。\par
日暮漢宮傳蠟燭，輕煙散入五侯家。\par
江南春·（唐）杜牧\par
千里鶯啼綠映紅，水村山郭酒旗風。\par
南朝四百八十寺，多少樓臺煙雨中。\par
上高侍郎·（唐）高蟾\par
天上碧桃和露種，日邊紅杏倚雲栽。\par
芙蓉生在秋江上，不向東風怨未開。\par
絕句·（宋）僧志安\par
古木陰中系短篷，杖藜扶我過橋東。\par
沾衣欲溼杏花雨，吹面不寒楊柳風。\par
遊小園不值·（宋）葉紹翁\par
應嫌屐齒印蒼苔，十扣柴扉九不開。\par
春色滿園關不住，一枝紅杏出牆來。\par
客中行·（唐）李白\par
蘭陵美酒鬱金香，玉碗盛來琥珀光。\par
但使主人能醉客，不知何處是他鄉。\par
題屏·（宋）劉季孫\par
呢喃燕子語樑間，底事來驚夢裏閒。\par
說與旁人渾不解，杖藜攜酒看芝山。\par
絕句漫興·（唐）杜甫\par
腸斷春江欲盡頭，杖藜徐步立芳洲。\par
顛狂柳絮隨風舞，輕薄桃花逐水流。\par
慶全庵桃花·（宋）謝枋得\par
尋得桃源好避秦，桃紅又是一年春。\par
花飛莫遣隨流水，怕有漁郎來問津。\par
玄都觀桃花·（唐）劉禹錫\par
紫陌紅塵拂面來，無人不道看花回。\par
玄都觀裏桃千樹，盡是劉郎去後栽。\par
再遊玄都觀·（唐）劉禹錫\par
百畝庭中半是苔，桃花淨盡菜花開。\par
種桃道士歸何處，前度劉郎今又來。\par
滁州西澗·（唐）韋應物\par
獨憐幽草澗邊生，上有黃鸝深樹鳴。\par
春潮帶雨晚來急，野渡無人舟自橫。\par
花影·（宋）蘇軾\par
重重疊疊上瑤臺，幾度呼童掃不開。\par
剛被太陽收拾去，卻教明月送將來。\par
北山·（宋）王安石\par
北山輸綠漲橫陂，直塹回塘灩灩時。\par
細數落花因坐久，緩尋芳草得歸遲。\par
湖上·（宋）徐元傑\par
花開紅樹亂鶯啼，草長平湖白鷺飛。\par
風日晴和人意好，夕陽簫鼓幾船歸。\par
漫興·（唐）杜甫\par
糝徑楊花鋪白氈，點溪荷葉疊青錢。\par
筍根稚子無人見，沙上鳧雛傍母眠。\par
春晴·（宋）王駕\par
雨前初見花間蕊，雨後全無葉底花。\par
蜂蝶紛紛過牆去，卻疑春色在鄰家。\par
春暮·（宋）曹豳\par
門外無人問落花，綠陰冉冉遍天涯。\par
林鶯啼到無聲處，青草池塘獨聽蛙。\par
落花·（宋）朱淑貞\par
連理枝頭花正開，妒花風雨便相催。\par
願教青帝常爲主，莫遣紛紛點翠苔。\par
春暮遊小園·（宋）王淇\par
一從梅粉褪殘妝，塗抹新紅上海棠。\par
開到荼蘼花事了，絲絲夭棘出莓牆。\par
鶯梭·（宋）劉克莊\par
擲柳遷喬太有情，交交時作弄機聲。\par
洛陽三月花如錦，多少工夫織得成。\par
暮春即事·（宋）葉採\par
雙雙瓦雀行書案，點點楊花入硯池。\par
閒坐小窗讀周易，不知春去幾多時。\par
登山·（唐）李涉\par
終日昏昏醉夢間，忽聞春盡強登山。\par
因過竹院逢僧話，又得浮生半日閒。\par
蠶婦吟·（唐）謝枋得\par
子規啼徹四更時，起視蠶稠怕葉稀。\par
不信樓頭楊柳月，玉人歌舞未曾歸。\par
晚春·（唐）韓愈\par
草木知春不久歸，百般紅紫鬥芳菲。\par
楊花榆莢無才思，惟解漫天作雪飛。\par
傷春·（宋）楊萬里\par
準擬今春樂事濃，依然枉卻一東風。\par
年年不帶看花眼，不是愁中即病中。\par
送春·（宋）王令\par
三月殘花落更開，小檐日日燕飛來。\par
子規夜半猶啼血，不信東風喚不回。\par
三月晦日送春·（唐）賈島\par
三月正當三十日，風光別我苦吟身。\par
共君今夜不須睡，未到曉鍾猶是春。\par
客中初夏·（宋）司馬光\par
四月清和雨乍晴，南山當戶轉分明。\par
更無柳絮因風起，惟有葵花向日傾。\par
有約·（宋）司馬光\par
黃梅時節家家雨，青草池塘處處蛙。\par
有約不來過夜半，閒敲棋子落燈花。\par
初夏睡起·（宋）楊萬里\par
梅子流酸濺齒牙，芭蕉分綠上窗紗。\par
日長睡起無情思，閒看兒童捉柳花。\par
三衢道中·（宋）曾幾\par
梅子黃時日日晴，小溪泛盡卻山行。\par
綠陰不減來時路，添得黃鸝四五聲。\par
即景·（宋）朱淑貞\par
竹搖清影罩幽窗，兩兩時禽噪夕陽。\par
謝卻海棠飛盡絮，困人天氣日初長。\par
夏日·（宋）戴敏\par
乳鴨池塘水淺深，熟梅天氣半晴陰。\par
東園載酒西園醉，摘盡枇杷一樹金。\par
晚樓閒坐·（宋）王安石\par
四顧山光接水光，憑欄十里芰荷香。\par
清風明月無人管，並作南來一味涼。\par
山居夏日·（唐）高駢\par
綠樹陰濃夏日長，樓臺倒影入池塘。\par
水晶簾動微風起，滿架薔薇一院香。\par
田家·（宋）范成大\par
晝出耘田夜績麻，村莊兒女各當家。\par
童孫未解供耕織，也傍桑陰學種瓜。\par
村居即事·（宋）范成大\par
綠遍山原白滿川，子規聲裏雨如煙。\par
鄉村四月閒人少，才了蠶桑又插田。\par
題榴花·（宋）朱熹\par
五月榴花照眼明，枝間時見子初成。\par
可憐此地無車馬，顛倒蒼苔落絳英。\par
村晚·（宋）雷震\par
草滿池塘水滿陂，山銜落日浸寒漪。\par
牧童歸去橫牛背，短笛無腔信口吹。\par
茅檐·（宋）王安石\par
茅檐常掃淨無苔，花木成蹊手自栽。\par
一水護田將綠繞，兩山排闥送青來。\par
烏衣巷·（唐）劉禹錫\par
朱雀橋邊野草花，烏衣巷口夕陽斜。\par
舊時王謝堂前燕，飛入尋常百姓家。\par
送元二使安西·（唐）王維\par
渭城朝雨浥輕塵，客舍青青柳色新。\par
勸君更盡一杯酒，西出陽關無故人。\par
與史朗中飲聽黃鶴樓上吹笛·（唐）李白\par
一爲遷客去長沙，西望長安不見家。\par
黃鶴樓中吹玉笛，江城五月落梅花。\par
題淮南寺·（宋）程顥\par
南去北來休便休，白蘋吹盡楚江秋。\par
道人不是悲秋客，一任晚山相對愁。\par
秋月·（宋）程顥\par
清溪流過碧山頭，空水澄鮮一色秋。\par
隔斷紅塵三十里，白雲紅葉兩悠悠。\par
七夕·（宋）楊樸\par
未會牽牛意若何，須邀織女弄金梭。\par
年年乞與人間巧，不道人間巧已多。\par
立秋·（宋）劉翰\par
乳鴉啼散玉屏空，一枕新涼一扇風。\par
睡起秋聲無覓處，滿階梧葉月明中。\par
秋夕·（唐）杜牧\par
銀燭秋光冷畫屏，輕羅小扇撲流螢。\par
天街夜色涼如水，臥看牽牛織女星。\par
中秋月·（唐）杜牧\par
暮雲收盡溢清寒，銀漢無聲轉玉盤。\par
此生此夜不長好，明月明年何處看？\par
江樓有感·（唐）趙嘏\par
獨上江樓思渺然，月光如水水如天。\par
同來玩月人何在，風景依稀似去年。\par
題臨安邸·（宋）林升\par
山外青山樓外樓，西湖歌舞幾時休。\par
暖風薰得遊人醉，直把杭州作汴州。\par
西湖·（宋）蘇軾\par
畢竟西湖六月中，風光不與四時同。\par
接天蓮葉無窮碧，映日荷花別樣紅。\par
飲湖上初晴後雨·（宋）蘇軾\par
水光瀲灩晴方好，山色空濛雨亦奇。\par
欲把西湖比西子，淡妝濃抹總相宜。\par
入直召對選德殿賜茶而退·（宋）周必大\par
綠槐夾道集昏鴉，敕使傳宣坐賜茶。\par
歸到玉堂清不寐，月鉤初上紫薇花。\par
夏日登車蓋亭·（宋）蔡確\par
紙屏石枕竹方牀，手倦拋書午夢長。\par
睡起莞然成獨笑，數聲漁笛在滄浪。\par
真玉堂作·（宋）洪諮夔\par
禁門深鎖寂無譁，濃墨淋漓兩相麻。\par
唱徹五更天未曉，一墀月浸紫薇花。\par
竹樓·（唐）李嘉佑\par
傲吏身閒笑五侯，西江取竹起高樓。\par
南風不用蒲葵扇，紗帽閒眠對水鷗。\par
直中書省·（唐）白居易\par
絲綸閣下文章靜，鐘鼓樓中刻漏長。\par
獨坐黃昏誰是伴，紫薇花對紫薇郎。\par
觀書有感（其一）·（宋）朱熹\par
半畝方塘一鑑開，天光雲影共徘徊。\par
問渠那得清如許，爲有源頭活水來。\par
泛舟·（宋）朱熹\par
昨夜江邊春水生，艨艟鉅艦一毛輕。\par
向來枉費推移力，此日中流自在行。\par
冷泉亭·（宋）林稹\par
一泓清可沁詩脾，冷暖年來只自知。\par
流出西湖載歌舞，回頭不似在山時。\par
冬景·（宋）蘇軾\par
荷盡已無擎雨蓋，菊殘猶有傲霜枝。\par
一年好景君須記，最是橙黃橘綠時。\par
楓橋夜泊·（唐）張繼\par
月落烏啼霜滿天，江楓漁火對愁眠。\par
姑蘇城外寒山寺，夜半鐘聲到客船。\par
寒夜·（宋）杜小山\par
寒夜客來茶當酒，竹爐湯沸火初紅。\par
尋常一樣窗前月，纔有梅花便不同。\par
霜月·（唐）李商隱\par
初聞徵雁已無蟬，百尺樓臺水接天。\par
青女素娥俱耐冷，月中霜裏鬥嬋娟。\par
梅·（宋）王淇\par
不受塵埃半點侵，竹籬茅舍自甘心。\par
只因誤識林和靖，惹得詩人說到今。\par
早春·（宋）白玉蟾\par
南枝才放兩三花，雪裏吟香弄粉些。\par
淡淡著煙濃著月，深深籠水淺籠沙。\par
雪梅·（宋）盧梅坡\par
梅雪爭春未肯降，騷人閣筆費評章。\par
梅須遜雪三分白，雪卻輸梅一段香。\par
有梅無雪不精神，有雪無詩俗了人。\par
日暮詩成天又雪，與梅井作十分春。\par
答鍾弱翁·（宋）牧童\par
草鋪橫野六七裏，笛弄晚風三四聲。\par
歸來飽飯黃昏後，不脫蓑衣臥月明。\par
秦淮夜泊·（唐）杜牧\par
煙籠寒水月籠沙，夜泊秦淮近酒家。\par
商女不知亡國恨，隔江猶唱後庭花。\par
歸雁·（唐）錢起\par
瀟湘何事等閒回，水碧沙明兩岸苔。\par
二十五絃彈夜月，不勝清怨卻飛來。\par
題壁·無名氏\par
一團茅草亂蓬蓬，驀地燒天驀地空。\par
爭似滿爐煨榾柮，漫騰騰地暖烘烘。\par

\section{第4则}
早朝大明宮·（唐）賈至\par
銀燭朝天紫陌長，禁城春色曉蒼蒼。\par
千條弱柳垂青鎖，百囀流鶯繞建章。\par
劍佩聲隨玉墀步，衣冠身惹御爐香。\par
共沐恩波鳳池上，朝朝染翰侍君王。\par
和賈舍人早朝·（唐）杜甫\par
五夜漏聲催曉箭，九重春色醉仙桃。\par
旌旗日暖龍蛇動，宮殿風微燕雀高。\par
朝罷香菸攜滿袖，詩成珠玉在揮毫。\par
欲知世掌絲綸美，池上於今有鳳毛。\par
和賈舍人早朝·（唐）王維\par
絳幘雞人報曉籌，尚衣方進翠雲裘。\par
九天閶闔開宮殿，萬國衣冠拜冕旒。\par
日色才臨仙掌動，香菸欲傍袞龍浮。\par
朝罷須裁五色詔，佩聲歸到鳳池頭。\par
和賈舍人早朝·（唐）岑參\par
雞鳴紫陌曙光寒，鶯囀皇州春色闌。\par
金闕曉鍾開萬戶，玉階仙仗擁千官。\par
花迎劍佩星初落，柳拂旌旗露未乾。\par
獨有鳳凰池上客，陽春一曲和皆難。\par
上元應制·（宋）蔡襄\par
高列千峯寶炬森，端門方喜翠華臨。\par
宸遊不爲三元夜，樂事還同萬衆心。\par
天上清光留此夕，人間和氣閣春陰。\par
要知盡慶華封祝，四十餘年惠愛深。\par
上元應制·（宋）王淇\par
雪消華月滿仙台，萬燭當樓寶扇開。\par
雙鳳雲中扶輦下，六鰲海上駕山來。\par
鎬京春酒沾周宴，汾水秋風陋漢才。\par
一曲昇平人盡樂，君王又進紫霞杯。\par
侍宴·（唐）沈佺期\par
皇家貴主好神仙，別業初開雲漢邊。\par
山出盡如鳴鳳嶺，池成不讓飲龍川。\par
妝樓翠幌教春住，舞閣金鋪借日懸。\par
敬從乘輿來此地，稱觴獻壽樂鈞天。\par
答丁元珍·（宋）歐陽修\par
春風疑不到天涯，二月山城未見花。\par
殘雪壓枝猶有橘，凍雷驚筍欲抽芽。\par
夜聞啼雁生鄉思，病入新年感物華。\par
曾是洛陽花下客，野芳雖晚不須嗟。\par
插花吟·（宋）邵雍\par
頭上花枝照酒卮，酒卮中有好花枝。\par
身經兩世太平日，眼見四朝全盛時。\par
況復筋骸粗康健，那堪時節正芳菲。\par
酒涵花影紅光溜，爭忍花前不醉歸。\par
寓意·（宋）晏殊\par
油壁香車不再逢，峽雲無跡任西東。\par
梨花院落溶溶月，柳絮池塘淡淡風。\par
幾日寂寥傷酒後，一番蕭瑟禁菸中。\par
魚書欲寄何由達，水遠山長處處同。\par
寒食書事·（宋）趙鼎\par
寂寂柴門村落裏，也教插柳紀年華。\par
禁菸不到粵人國，上冢亦攜龐老家。\par
漢寢唐陵無麥飯，山溪野徑有梨花。\par
一樽竟藉青苔臥，莫管城頭奏暮笳。\par
清明·（宋）黃庭堅\par
佳節清明桃李笑，野田荒冢只生愁。\par
雷驚天地龍蛇蟄，雨足郊原草木柔。\par
人乞祭餘驕妾婦，士甘焚死不公侯。\par
賢愚千載知誰是，滿眼蓬蒿共一丘。\par
清明·（宋）高菊卿\par
南北山頭多墓田，清明祭掃各紛然。\par
紙灰飛作白蝴蝶，淚血染成紅杜鵑。\par
日落狐狸眠冢上，夜歸兒女笑燈前。\par
人生有酒須當醉，一滴何曾到九泉。\par
郊行即事·（宋）程顥\par
芳原綠野恣行時，春入遙山碧四圍。\par
興逐亂紅穿柳巷，困臨流水坐苔磯。\par
莫辭盞酒十分勸，只恐風花一片飛。\par
況是清明好天氣，不妨遊衍莫忘歸。\par
鞦韆·（宋）僧惠洪\par
畫架雙裁翠絡偏，佳人春戲小樓前。\par
飄揚血色裙拖地，斷送玉容人上天。\par
花皮潤沾紅杏雨，彩繩斜掛綠楊煙。\par
下來閒處從容立，疑是蟾宮謫神仙。\par
曲江對酒·（唐）杜甫\par
一片花飛減卻春，風飄萬點正愁人。\par
且看欲盡花經眼，莫厭傷多酒入脣。\par
江上小堂巢翡翠，苑邊高冢臥麒麟。\par
細推物理須行樂，何用浮名絆此身。\par
朝回日日典春衣，每日江頭盡醉歸。\par
酒債尋常行處有，人生七十古來稀。\par
穿花蛺蝶深深見，點水蜻蜓款款飛。\par
傳與風光共流轉，暫時相賞莫相違。\par
黃鶴樓·（唐）崔顥\par
昔人已乘黃鶴去，此地空餘黃鶴樓。\par
黃鶴一去不復返，白雲千載空悠悠。\par
晴川歷歷漢陽樹，芳草萋萋鸚鵡州。\par
日暮鄉關何處是，煙波江上使人愁。\par
旅懷·（唐）崔塗\par
水流花謝兩無情，送盡東風過楚城。\par
蝴蝶夢中家萬里，杜鵑枝上月三更。\par
故園書動經年絕，華髮春催兩鬢生。\par
自是不歸歸便得，五湖煙景有誰爭。\par
答李儋·（唐）韋應物\par
去年花裏逢君別，今日花開又一年。\par
世事茫茫難自料，春愁黯黯獨成眠。\par
身多疾病思田裏，邑有流亡愧俸錢。\par
聞道欲來相問訊，西樓望月幾回圓。\par
江村·（唐）杜甫\par
清江一曲抱村流，長夏江村事事幽。\par
自去自來樑上燕，相親相近水中鷗。\par
老妻畫紙爲棋局，稚子敲針作釣鉤。\par
多病所須惟藥物，微軀此外更何求。\par
夏日·（宋）張文潛\par
長夏江村風日清，檐牙燕雀已生成。\par
蝶衣曬粉花枝午，蛛網添絲屋角晴。\par
落落疏廉邀月影，嘈嘈虛枕納溪聲。\par
久斑兩鬢如霜雪，直欲樵漁過此生。\par
積雨輞川莊作·（唐）王維\par
積雨空林煙火遲，蒸藜炊黍餉東菑。\par
漠漠水田飛白鷺，陰陰夏木囀黃鸝。\par
山中習靜合朝槿，松下清齋折露葵。\par
野老與人爭席罷，海鷗何事更相疑。\par
新竹·（宋）陸游\par
插棘編籬謹護持，養成寒碧映漣漪。\par
清風掠地秋先到，赤日行天午不知。\par
解籜時聞聲簌簌，放梢初見影離離。\par
歸閒我欲頻來此，枕簟仍教到處隨。\par
偶成·（宋）程顥\par
閒來無事不從容，睡覺東窗日已紅。\par
萬物靜觀皆自得，四時佳興與人同。\par
道通天地有形外，思入風雲變態中。\par
富貴不淫貧賤樂，男兒到此是豪雄。\par
表兄話舊·（唐）竇叔向\par
夜合花開香滿庭，夜深微雨醉初醒。\par
遠書珍重何由達，舊事淒涼不可聽。\par
去日兒童皆長大，昔年親友半凋零。\par
明朝又是孤舟別，愁見河橋酒幔青。\par
遊月陂·（宋）程顥\par
月坡堤上四徘徊，北有中天百尺臺。\par
萬物已隨秋氣改，一樽聊爲晚涼開。\par
水心雲影閒相照，林下泉聲靜自來。\par
世事無端何足計，但逢佳節約重陪。\par
秋興八首·（唐）杜甫\par
千家山郭靜朝暉，日日江樓坐翠微。\par
信宿漁人還泛泛，清秋燕子故飛飛。\par
匡衡抗疏功名薄，劉向傳經心事違。\par
同學少年多不賤，五陵裘馬自輕肥。\par
蓬萊宮闕對南山，承露金莖霄漢間。\par
西望瑤池降王母，東來紫氣滿函關。\par
雲移雉尾開宮扇，日繞龍鱗識聖顏。\par
一臥滄江驚歲晚，幾回青瑣點朝班。\par
玉露凋傷楓樹林，巫山巫峽氣蕭森。\par
江間波浪兼天涌，塞上風雲接地陰。\par
叢菊兩開他日淚，孤舟一系故園心。\par
寒衣處處催刀尺，白帝城高急暮砧。\par
昆明池水漢時功，武帝旌旗在眼中。\par
織女機絲虛夜月，石鯨鱗甲動秋風。\par
波飄菰米沉雲黑，露冷蓮房墜粉紅。\par
關塞極天惟鳥道，江湖滿地一漁翁。\par
月夜舟中·（宋）戴復古\par
滿船明月浸虛空，綠水無痕夜氣衝。\par
詩思浮沈檣影裏，夢魂搖拽櫓聲中。\par
星辰冷落碧潭水，鴻雁悲鳴紅蓼風。\par
數點漁燈依古岸，斷橋垂露滴梧桐。\par
長安秋望·（唐）趙嘏\par
雲物淒涼拂署流，漢家宮闕動高秋。\par
殘星幾點雁橫塞，長笛一聲人倚樓。\par
紫豔半開籬菊靜，紅衣落盡渚蓮愁。\par
鱸魚正美不歸去，空戴南冠學楚囚。\par
新秋·（唐）杜甫\par
火雲猶未斂奇峯，欹枕初驚一葉風。\par
幾處園林蕭瑟裏，誰家砧杵寂寥中。\par
蟬聲斷續悲殘月，螢焰高低照暮空。\par
賦就金門期再獻，夜深搔首嘆飛蓬。\par
中秋·（宋）李樸\par
皓魄當空寶鏡升，雲間仙籟寂無聲。\par
平分秋色一輪滿，長伴雲衢千里明。\par
狡兔空從弦外落，妖蟆休向眼前生。\par
靈槎擬約同攜手，更待銀河澈底清。\par
九日藍耕會飲·（唐）杜甫\par
老去悲秋強自寬，興來今日盡君歡。\par
羞將短髮還吹帽，笑倩旁人爲正冠。\par
藍水遠從千澗落，玉山高並兩峯寒。\par
明年此會知誰健，醉把茱萸仔細看。\par
秋思·（宋）陸游\par
利慾驅人萬火牛，江湖浪跡一沙鷗。\par
日長似歲閒方覺，事大如天醉亦休。\par
砧杵敲殘深巷月，梧桐搖落故園秋。\par
欲舒老眼無高處，安得元龍百尺樓。\par
與朱山人·（唐）杜甫\par
錦裏先生烏角巾，園收芋慄未全貧。\par
慣看賓客兒童喜，得食階除鳥雀馴。\par
秋水才深四五尺，野航恰受兩三人。\par
白沙翠竹江村暮，相送柴門月色新。\par
聞笛·（唐）趙嘏\par
誰家吹笛畫樓中，斷續聲隨斷續風。\par
響遏行雲橫碧落，清和冷月到簾櫳。\par
興來三弄有桓子，賦就一篇懷馬融。\par
曲罷不知人在否，餘音嘹亮尚飄空。\par
冬景·（宋）劉克莊\par
晴窗早覺愛朝曦，竹外秋聲漸作威。\par
命僕安排新暖閣，呼童熨貼舊寒衣。\par
葉浮嫩綠酒初熟，橙切香黃蟹正肥。\par
蓉菊滿園皆可羨，賞心從此莫相違。\par
冬至·（唐）杜甫\par
天時人事日相催，冬至陽生春又來。\par
刺繡五紋添弱線，吹葭六管動飛灰。\par
岸容待臘將舒柳，山意衝寒欲放梅。\par
雲物不殊鄉國異，教兒且覆掌中杯。\par
梅花·（宋）林逋\par
衆芳搖落獨鮮妍，占斷風情向小園。\par
疏影橫斜水清淺，暗香浮動月黃昏。\par
霜禽欲下先偷眼，粉蝶如知合斷魂。\par
幸有微吟可相狎，不須檀板共金樽。\par
左遷至藍關示侄孫湘·（唐）韓愈\par
一封朝奏九重天，夕貶潮陽路八千。\par
本爲聖朝除弊政，敢將衰朽惜殘年。\par
雲橫秦嶺家何在，雪擁藍關馬不前。\par
知汝遠來應有意，好收吾骨瘴江邊。\par
干戈·（宋）王中\par
干戈未定欲何之，一事無成兩鬢絲。\par
蹤跡大綱王粲傳，情懷小樣杜陵詩。\par
鶺鴒音斷人千里，烏鵲巢寒月一枝。\par
安得中山千日酒，酩然直到太平時。\par
歸隱·（宋）陳摶\par
十年蹤跡走紅塵，回首青山入夢頻。\par
紫綬縱榮爭及睡，朱門雖富不如貧。\par
愁聞劍戟扶危主，悶聽笙歌聒醉人。\par
攜取舊書歸舊隱，野花啼鳥一般春。\par
山中寡婦·（唐）杜荀鶴\par
夫因兵亂守蓬茅，麻苧裙衫鬢髮焦。\par
桑柘廢來猶納稅，田園荒盡尚徵苗。\par
時挑野菜和根煮，旋砍生柴帶葉燒。\par
任是深山最深處，也應無計避徵徭。\par
送天師·（明）朱權\par
霜落芝城柳影疏，殷勤送客出鄱湖。\par
黃金甲鎖雷霆印，紅錦韜纏日月符。\par
天上曉行騎只鶴，人間夜宿解雙鳧。\par
匆匆歸到神仙府，爲問蟠桃熟也無。\par
送毛伯溫·（明）明世宗\par
大將南征膽氣豪，腰橫秋水雁翎刀。\par
風吹鼉鼓山河動，電閃旌旗日月高。\par
天上麒麟原有種，穴中螻蟻豈能逃。\par
太平待詔歸來日，朕與先生解戰袍。\par

\end{document}